\chapter{Python 代码示例}

本附录给出若干可直接运行的 Python 示例。默认读者已安装 \texttt{numpy}、\texttt{matplotlib}，其中第一题还需要 \texttt{scipy}。

\section{用 \texttt{scipy.integrate.odeint} 求解阻尼振子}

考虑方程
\[
 m\ddot x+b\dot x+kx=0.
\]
将其化为一阶方程组后，可用 \texttt{odeint} 直接求解。

\begin{lstlisting}[language=Python]
import numpy as np
import matplotlib.pyplot as plt
from scipy.integrate import odeint

m = 1.0
b = 0.4
k = 4.0


def rhs(state, t, m, b, k):
    x, v = state
    dxdt = v
    dvdt = -(b / m) * v - (k / m) * x
    return [dxdt, dvdt]


t = np.linspace(0.0, 20.0, 2000)
y0 = [1.0, 0.0]
sol = odeint(rhs, y0, t, args=(m, b, k))
x = sol[:, 0]
v = sol[:, 1]

plt.figure(figsize=(8, 4))
plt.plot(t, x, label='x(t)')
plt.plot(t, v, label='v(t)', alpha=0.8)
plt.xlabel('t')
plt.ylabel('state')
plt.title('Damped Oscillator solved by odeint')
plt.grid(True)
plt.legend()
plt.tight_layout()
plt.show()
\end{lstlisting}

\section{用 RK4 手写求解抛体加空气阻力}

这里采用速度平方阻力模型，阻力方向与速度相反：
\[
\vb{F}_d=-c\abs{\vb{v}}\vb{v}.
\]

\begin{lstlisting}[language=Python]
import numpy as np
import matplotlib.pyplot as plt


g = 9.81
c_drag = 0.02
m = 1.0
h = 0.01
T = 12.0

speed0 = 50.0
angle = np.deg2rad(45.0)
state = np.array([0.0, 0.0,
                  speed0 * np.cos(angle),
                  speed0 * np.sin(angle)], dtype=float)


def rhs(s):
    x, y, vx, vy = s
    v = np.hypot(vx, vy)
    ax = -(c_drag / m) * v * vx
    ay = -g - (c_drag / m) * v * vy
    return np.array([vx, vy, ax, ay], dtype=float)


def rk4_step(s, h):
    k1 = rhs(s)
    k2 = rhs(s + 0.5 * h * k1)
    k3 = rhs(s + 0.5 * h * k2)
    k4 = rhs(s + h * k3)
    return s + (h / 6.0) * (k1 + 2*k2 + 2*k3 + k4)


ts = [0.0]
xs = [state[0]]
ys = [state[1]]

while ts[-1] < T and state[1] >= 0.0:
    state = rk4_step(state, h)
    ts.append(ts[-1] + h)
    xs.append(state[0])
    ys.append(state[1])

plt.figure(figsize=(7, 4))
plt.plot(xs, ys, label='with drag')
plt.xlabel('x')
plt.ylabel('y')
plt.title('Projectile with Air Drag (RK4)')
plt.grid(True)
plt.legend()
plt.tight_layout()
plt.show()
\end{lstlisting}

\section{大角度单摆的数值解与小角近似对比}

完整方程为
\[
\ddot \theta+\frac{g}{L}\sin\theta=0,
\]
小角近似则将 $\sin\theta$ 替换为 $\theta$。

\begin{lstlisting}[language=Python]
import numpy as np
import matplotlib.pyplot as plt


g = 9.81
L = 1.0
h = 0.01
T = 12.0
state = np.array([1.2, 0.0], dtype=float)  # theta, omega


def rhs_exact(s):
    theta, omega = s
    return np.array([omega, -(g / L) * np.sin(theta)], dtype=float)


def rhs_small(s):
    theta, omega = s
    return np.array([omega, -(g / L) * theta], dtype=float)


def rk4_step(rhs, s, h):
    k1 = rhs(s)
    k2 = rhs(s + 0.5 * h * k1)
    k3 = rhs(s + 0.5 * h * k2)
    k4 = rhs(s + h * k3)
    return s + (h / 6.0) * (k1 + 2*k2 + 2*k3 + k4)


t_vals = np.arange(0.0, T + h, h)
state_exact = state.copy()
state_small = state.copy()
exact = []
small = []

for _ in t_vals:
    exact.append(state_exact[0])
    small.append(state_small[0])
    state_exact = rk4_step(rhs_exact, state_exact, h)
    state_small = rk4_step(rhs_small, state_small, h)

plt.figure(figsize=(8, 4))
plt.plot(t_vals, exact, label='nonlinear pendulum')
plt.plot(t_vals, small, '--', label='small-angle approximation')
plt.xlabel('t')
plt.ylabel('theta')
plt.title('Pendulum: Exact vs Small-Angle Approximation')
plt.grid(True)
plt.legend()
plt.tight_layout()
plt.show()
\end{lstlisting}

\section{RLC 电路暂态响应}

对串联 RLC 电路，电荷满足
\[
 L\ddot q+R\dot q+\frac{1}{C}q=0.
\]
下例展示欠阻尼情形的暂态衰减振荡。

\begin{lstlisting}[language=Python]
import numpy as np
import matplotlib.pyplot as plt


L = 0.5
R = 1.0
C = 0.2
h = 0.002
T = 10.0
state = np.array([1.0, 0.0], dtype=float)  # q, i


def rhs(s):
    q, i = s
    dqdt = i
    didt = -(R / L) * i - q / (L * C)
    return np.array([dqdt, didt], dtype=float)


def rk4_step(s, h):
    k1 = rhs(s)
    k2 = rhs(s + 0.5 * h * k1)
    k3 = rhs(s + 0.5 * h * k2)
    k4 = rhs(s + h * k3)
    return s + (h / 6.0) * (k1 + 2*k2 + 2*k3 + k4)


t_vals = np.arange(0.0, T + h, h)
q_vals = []
i_vals = []

for _ in t_vals:
    q_vals.append(state[0])
    i_vals.append(state[1])
    state = rk4_step(state, h)

plt.figure(figsize=(8, 4))
plt.plot(t_vals, q_vals, label='charge q(t)')
plt.plot(t_vals, i_vals, label='current i(t)', alpha=0.8)
plt.xlabel('t')
plt.ylabel('response')
plt.title('Transient Response of a Series RLC Circuit')
plt.grid(True)
plt.legend()
plt.tight_layout()
plt.show()
\end{lstlisting}

\section{带电粒子在磁场中的螺旋运动 3D 轨迹}

令均匀磁场为 $\vb{B}=B\vu{z}$，忽略电场，则洛伦兹力为
\[
 m\dv{\vb{v}}{t}=q\vb{v}\times\vb{B}.
\]
当速度具有平行与垂直于磁场的分量时，轨迹是螺旋线。

\begin{lstlisting}[language=Python]
import numpy as np
import matplotlib.pyplot as plt
from mpl_toolkits.mplot3d import Axes3D


q = 1.0
m = 1.0
Bz = 1.5
h = 0.01
T = 30.0
state = np.array([0.0, 0.0, 0.0,
                  1.0, 0.5, 0.8], dtype=float)  # x, y, z, vx, vy, vz


def rhs(s):
    x, y, z, vx, vy, vz = s
    ax = (q / m) * vy * Bz
    ay = -(q / m) * vx * Bz
    az = 0.0
    return np.array([vx, vy, vz, ax, ay, az], dtype=float)


def rk4_step(s, h):
    k1 = rhs(s)
    k2 = rhs(s + 0.5 * h * k1)
    k3 = rhs(s + 0.5 * h * k2)
    k4 = rhs(s + h * k3)
    return s + (h / 6.0) * (k1 + 2*k2 + 2*k3 + k4)


t_vals = np.arange(0.0, T + h, h)
xs, ys, zs = [], [], []

for _ in t_vals:
    xs.append(state[0])
    ys.append(state[1])
    zs.append(state[2])
    state = rk4_step(state, h)

fig = plt.figure(figsize=(7, 5))
ax = fig.add_subplot(111, projection='3d')
ax.plot(xs, ys, zs, lw=1.5)
ax.set_xlabel('x')
ax.set_ylabel('y')
ax.set_zlabel('z')
ax.set_title('Helical Motion in a Uniform Magnetic Field')
plt.tight_layout()
plt.show()
\end{lstlisting}

以上代码均可直接保存为 \texttt{.py} 文件运行。若希望提高可重复性，可进一步加入参数封装、命令行输入与数据导出功能。
